\documentclass[book.tex]{subfiles}

\begin{document}

\chapter{Universal Approximation Theorem}
\label{chap:univ_approx}
\noindent\rule{11cm}{0.4pt} \\
\textbf{Prerequisites:}  \autoref{chap:ml_basics}, \autoref{chap:grad_descent}  \\
\textbf{Difficulty Level:} *** \\
\noindent\rule{11cm}{0.4pt}
\vspace{5mm}

The universal approximation theorem is one of the most foundational theoretical results in deep learning. It proves that arbitrary functions can be approximated by fully connected deep networks. There are a number of extensions of the universal approximation theorem to other families of neural networks that are also useful.

\section{The Universal Approximation Theorem for Fully Connected Networks}

Start with the one layer network. Let $C(X, Y)$ denote the set of continuous functions from $X$ to $Y$, where $X$ is a compact space. Then it has been shown that sums of the form

\begin{align}
\sum_{j=1}^N a_j\sigma(w_jx + b_j)
\end{align}
can approximate any function in $C(X, Y)$ arbitrarily well for suitable choices of parameters. The proof of the theorem draws upon the fact that Fourier series offer good approximations to functions in $C(X, Y)$

\section{The Universal Approximation Theorem for Convolutional Networks}

Convolutional networks form a subset of fully connected networks with a special Toeplitz matrix structure for the weight matrix. A priori, it is not necessarily obvious a universal approximation theorem would hold for such systems, but this has been proven to be the case.

\section{The Universal Approximation Theorem for Recurrent Neural Networks}

Any open dynamic system can be approximated with arbitrary accuracy by a recurrent neural network. Here an open dynamic system is given by a set of update equations

\begin{align}
s_{t+1} &= g(s_t, x_t)\\
y_t &= h(s_t)
\end{align}

where $x_t$ is the input at time $t$, $s_t$ is the state at time $t$ and $y_t$ is the output at time $t$.

\section{The Universal Approximation Theorem for Graph Networks}

Function approximation on graphs requires a slightly more complicated framework to specify precisely. Two different graphs can be isomorphic; that is there can exist a map
\begin{align}
\phi: G \to H
\end{align}
between graphs $G$ and $H$ that preserves edges and labels on vertices. The universal approximation theorem for graphs states that any function that has the same value on isomorphic graphs can be approximated by a suitable graph neural network.
%Graph convolutional networks are as expressive as the Weisfeiler-Leman graph isomorphism test. 

%\textcolor{red}{TODO}

\end{document}